% Generated draft: requires mathematical and source review.
\documentclass[11pt]{article}
\usepackage[margin=25mm]{geometry}
\usepackage{fontspec}
\setmainfont{Latin Modern Roman}
\usepackage{amsmath,amssymb,mathtools,unicode-math}
\setmathfont{Latin Modern Math}
\usepackage{enumitem,needspace,xurl,hyperref}
\setlength{\parindent}{0pt}
\setlength{\parskip}{5pt}
\setlength{\emergencystretch}{3em}
\newcommand{\N}{\mathbb{N}}
\newcommand{\Z}{\mathbb{Z}}
\newcommand{\Q}{\mathbb{Q}}
\newcommand{\R}{\mathbb{R}}
\newcommand{\C}{\mathbb{C}}
\newcommand{\F}{\mathbb{F}}
\newcommand{\A}{\mathbb{A}}
\newcommand{\PP}{\mathbb{P}}
\newcommand{\E}{\mathbb{E}}
\newcommand{\eps}{\varepsilon}
\newcommand{\dd}{\,\mathrm{d}}
\begin{document}
\section*{UnsolvedMath ID: 700576}
\textbf{Source problem code:} catalogue entry
\subsection{Definitions and mathematical statement}
A finite projective plane is an incidence structure \(\Pi=(P,L,I)\), with points \(P\), lines \(L\), and incidence relation \(I\subseteq P\times L\), such that any two distinct points lie on exactly one common line, any two distinct lines meet in exactly one point, and there exist four points no three of which are collinear. Its order is \(q\) if every line contains exactly \(q+1\) points and every point lies on exactly \(q+1\) lines. A plane is Desarguesian if it is isomorphic to the projective plane over a division ring; for a finite field \(\F_q\), the corresponding plane is denoted by \(\operatorname{PG}(2,q)\).

The following is a qualified transcription of N. D. Podufalov's problem, with its unresolved construction data retained explicitly.

Part (b), the remaining question: The source refers to a group \(G\) in Podufalov's 1984 symposium abstract and to a group \(S\). It states that a projective plane of order \(3\) can be constructed by taking as lines the members of any class of subgroups of order \(2\cdot5\cdot11\cdot17\) conjugate under \(S\), and adjoining four further lines ``in a natural way.'' It also states that a similar construction gives projective planes of order \(p^n\) for every prime \(p\) and every natural number \(n\). Can this method be adapted to construct new projective planes?

Here the supplied record does not define \(G\), \(S\), their relationship, the point set, the incidence relation, the four additional lines, or the general construction for \(p^n\). It also does not specify the meaning of ``new,'' including the comparison class of previously known planes. Consequently, this question cannot be converted from the supplied record into a fully specified existence statement. The convention for natural numbers in part (b) is unspecified; inclusion of \(n=0\) would require clarification because order \(1\) is incompatible with the nondegenerate projective-plane definition above.

Part (a), the archived question: The source refers to the group \(G\) described in the archived solution of Problem 8.73. It states that a projective plane of order \(3\) is obtained by taking as lines the members of any conjugacy class of subgroups of order \(2\cdot5\cdot7\cdot11\), together with four lines added ``in a natural way.'' It further states that an analogous construction gives a projective plane of order \(p^n\) for every prime \(p\) and every positive integer \(n\). Does the resulting projective plane have the Galois property?

The supplied record does not define the Galois property or give the construction needed to determine the family quantified over in part (a). The different subgroup orders and references in parts (a) and (b) are preserved as supplied; no identification of their groups or constructions is assumed.

The supplied Archive of Solved Problems entry reports that the answer to part (a) is affirmative for \(n=1\), but is not always affirmative for \(n>1\). It further reports that, if a near-field of \(p^n\) elements exists, then some of the indicated planes of order \(p^n\) are non-Desarguesian. This is an attributed archive report, not an independently verified result here. The record does not specify whether ``near-field'' in that conditional assertion requires a near-field that is not a field, nor does it explicitly define the relationship between the Galois property and being Desarguesian.
\subsection{Short English statement}
Podufalov asks whether the indicated group construction can be adapted to produce new finite projective planes. The archived part asks whether its resulting planes have the Galois property and includes a supplied negative answer in general, with an affirmative answer for \(n=1\). Essential construction data, the Galois property, and the meaning of ``new'' are absent, so a complete formal statement requires further source material.
\subsection{Sources}
\begin{itemize}
\item Supplied local UnsolvedMath registry record: permanent ID \(700576\), source code \texttt{catalogue entry}, title ``Kourovka Notebook Problem 9.43.'' Proposed by N. D. Podufalov in 1984. The supplied registry labels the record \texttt{partially\_solved}; this label is not an independent mathematical assessment.
\item The Kourovka Notebook (Unsolved Problems in Group Theory), 9th issue (1984), Problem 9.43; the supplied record identifies an update in the 21st issue (2026), \texttt{arXiv:1401.0300v48}. These references and their contents have not been independently checked here.
\item N. D. Podufalov, Abstracts of the 9th All--Union Symposium on Group Theory, Moscow, 1984, pp. 113--114 (Russian). The supplied part (b) refers to this abstract for its group \(G\).
\item Archived solution of Kourovka Notebook Problem 8.73, cited in the supplied part (a) for its group \(G\); the solution itself is not included in the record.
\item Archive of Solved Problems entry for Problem 9.43, attributed to N. D. Podufalov, Letter of November 24, 1986. The supplied entry reports the qualification for \(n=1\), the negative answer in general, and the conditional near-field assertion.
\item Supplied status review checked on \texttt{2026-10-16}: it records solved part(s) in the archive and remaining part(s) in the main section, with evidence labelled ``published.'' This is supplied dated curation, not independently verified source or literature assessment.
\end{itemize}
\end{document}
